%------------------------------------------------------------------------------------
%  SECTION 8.2 : GESTION BIBLIOGRAPHIQUE AVEC BIBLATEX ET STYLE VANCOUVER
%------------------------------------------------------------------------------------
\section{Gestion Bibliographique Moderne et Style Vancouver}
\label{sec:bibliographie_vancouver}

La qualité des sources et leur référencement rigoureux constituent le socle de toute démarche scientifique.
L'ancien moteur \textsf{Bib\TeX{}} (1985), limité aux caractères ASCII 7 bits et dépendant de fichiers de style rigides (\texttt{.bst}), est aujourd'hui totalement dépassé.
\texttt{ClasseGuillaume} s'appuie sur le couple de référence contemporain : le paquet \textsf{biblatex} propulsé par le moteur de tri Unicode \textbf{Biber}.

\subsection{Indépendance absolue entre données brutes (\texttt{.bib}) et style de rendu}
\label{subsec:independance_bib_style}

L'un des atouts fondamentaux de l'écosystème \LaTeX{} réside dans la séparation étanche entre le fond et la forme :
\begin{itemize}
  \item Le fichier \texttt{bib/references.bib} constitue une \textbf{base de données brute et universelle}.
  Il est totalement \emph{agnostique} du style de publication et ne stocke que les métadonnées sémantiques (auteurs, titres, volumes, éditeurs, dates, identifiants DOI et URL).
  \item Le \textbf{style de restitution visuelle} (Vancouver, IEEE, APA, Chicago, Nature, etc.) est quant à lui gouverné de manière unique dans le document \LaTeX{} (option \texttt{style=vancouver} de \textsf{biblatex} dans la classe ou le préambule).
\end{itemize}

Ainsi, pour soumettre un même mémoire ou article à une revue exigeant le style auteur-date (ex. APA) plutôt qu'un style numérique séquentiel (Vancouver ou IEEE), il suffit de modifier un seul mot dans la configuration de \textsf{biblatex}, sans jamais devoir réécrire ou modifier la moindre ligne du fichier \texttt{references.bib}.

\subsection{Pourquoi le style Vancouver médical et technique ?}
\label{subsec:pourquoi_vancouver}

Dans les disciplines biomédicales, les sciences appliquées et l'ingénierie, le style de citation séquentiel numéroté dit \textbf{Style Vancouver} est la norme internationale préconisée par l'\emph{International Committee of Medical Journal Editors} (ICMJE).
Contrairement aux styles auteur-date (type APA ou Harvard) qui alourdissent les phrases par de longues parenthèses, le style Vancouver insère un numéro de référence sobre et compact entre crochets, sans interrompre la lecture.

\subsection{Diversité des modes de citation en contexte}
\label{subsec:modes_citations_contexte}

Selon la tournure syntaxique de la phrase, plusieurs macros de citation sont mises en œuvre :

\begin{itemize}
  \item \textbf{Citation narrative en position de sujet (\texttt{\textbackslash textcite}) :}
  Comme l'a démontré \textcite{knuth1984texbook} dans ses traités fondateurs, la précision typographique relève d'une exigence algorithmique.
  De même, \textcite{lamport1994latex} puis \textcite{mittelbach2004companion} ont démocratisé la séparation entre forme et contenu.
  
  \item \textbf{Citation factuelle globale (\texttt{\textbackslash cite} ou \texttt{\textbackslash parencite}) :}
  L'agencement des proportions de la page découle des règles de l'art typographique classique~\cite{bringhurst2004elements}.
  La suite KOMA-Script apporte une flexibilité inégalée pour les ouvrages complexes~\cite{kohm2024koma},
  tandis que \textsf{fontspec} régit le choix des polices OpenType sous Lua\LaTeX{}~\cite{robertson2020fontspec}.
  
  \item \textbf{Études médicales et essais cliniques :}
  Dans le domaine de la santé, le diagnostic des pathologies gynécologiques s'appuie sur des méta-analyses rigoureuses~\cite{kanti2024ultrasound} et l'évaluation prospective de biomarqueurs sanguins~\cite{burghaus2024}.
  Par ailleurs, l'impact des directives cliniques sur les cohortes observationnelles a été modélisé récemment~\cite{nieva2024impact}.

  \item \textbf{Thèses de doctorat et manuels de référence :}
  Les problématiques d'optimisation documentaire sous Lua\LaTeX{} sont approfondies dans la thèse de doctorat de \textcite{winant2026thesis}.
  Le manuel officiel de référence Lua\TeX{} détaille quant à lui l'ensemble des primitives de bas niveau~\cite{luatex2024manual}.
\end{itemize}

\subsection{Francisation soignée et liens de retour \texttt{backref}}
\label{subsec:francisation_backref}

La classe personnalise le moteur \textsf{biblatex} pour une intégration francophone sans faille :
\begin{itemize}
  \item Les dates de consultation des ressources numériques affichent automatiquement « [consulté le \dots] » à partir du champ \texttt{urldate}.
  \item Les adresses web affichent la mention « disponible sur : » suivie du lien cliquable césuré proprement par \textsf{xurl}.
  \item \textbf{La fonctionnalité majeure \texttt{backref = true} :}
  Dans la section bibliographique en fin de document, chaque référence comporte automatiquement \textbf{la liste cliquable des numéros de pages} où elle a été citée dans le corps de l'ouvrage !
  Cette fonctionnalité permet aux membres du jury de vérifier instantanément dans quel contexte précis chaque source a été mobilisée.
\end{itemize}

Les citations de texte s'appuient sur le paquet \textsf{csquotes} avec la commande \texttt{\textbackslash enquote\{...\}}, produisant automatiquement des guillemets français dits en chevrons («~\dots~») avec espaces insécables normées.

